\section{Conclusion}

This work has demonstrated that the different coverage planners each have their strengths and weaknesses: The morphology-based skeleton coverage planner excelled in cluttered environments with tight spaces, but the resulting path did not provide complete coverage of the free space in larger rooms. conversely, the spanning tree coverage planner achieved better coverage at the expense a large number of sharp turns, reducing localisation accuracy.

The developed perception pipeline successfully localised and classified small objects, demonstrating that this combination is an effective solution for object recognition for the MIRTE Master.

The modifications to the manipulator have increased its accuracy nearly two orders of magnitude, as the system no longer relies on approximate joint values. This together with the improved gripper surface made it so objects are grabbable even at odd angles.

Overall, this work has successfully demonstrated the MIRTE Master can accommodate all components necessary to realise a laboratory-cleaning robot, and is therefore a viable platform for performing this task.
